\chapter{Épreuve chronométrée 2}\label{ch:iv1:timed-test-2}

\begin{ivtrains}
Faire l'épreuve d'une seule traite, sans notes ni calculatrice.
\end{ivtrains}

\ivtestsection{A}{6}{40}

\begin{question}\label{iq:iv1:timed-test-2:1}
Une urne contient 5 boules rouges et 3 boules bleues. Vous tirez sans remise jusqu'à ce qu'une des deux couleurs soit épuisée. En moyenne, combien de boules restent-elles dans l'urne lorsque vous vous arrêtez~?
\end{question}

\begin{question}\label{iq:iv1:timed-test-2:2}
Vous lancez à répétition une paire de dés. En moyenne, combien de lancers faut-il pour observer à la fois une somme de 7 et une somme de 11, dans n'importe quel ordre~?
\end{question}

\begin{question}\label{iq:iv1:timed-test-2:3}
Une action vaut 100 et vaudra 120 ou 90 à la période suivante. Les taux d'intérêt sont nuls. Quel est le prix d'un call de strike 100, et combien d'actions couvrent un call~?
\end{question}

\begin{question}\label{iq:iv1:timed-test-2:4}
$X$ et $Y$ sont indépendantes et uniformes sur $[0, 1]$. Quelle est la probabilité que $\sqrt{X} + \sqrt{Y} < 1$~?
\end{question}

\begin{question}\label{iq:iv1:timed-test-2:5}
Une pièce équilibrée est lancée 20 fois et donne exactement 10 piles. Quelle est la probabilité que exactement 5 de ces piles soient sorties dans les 10 premiers lancers~? Donnez une expression exacte et une valeur approchée.
\end{question}

\begin{question}\label{iq:iv1:timed-test-2:6}
Une pièce équilibrée est lancée jusqu'à obtenir la première pile. On vous indique que le nombre de lancers était pair. Sachant cela, quelle est l'espérance du nombre de lancers~?
\end{question}

\ivtestsection{B}{6}{60}

\begin{question}\label{iq:iv1:timed-test-2:7}
Un jeu de 52 cartes est distribué à quatre joueurs, 13 cartes chacun. Vous regardez votre main et voyez exactement deux as. Quelle est la probabilité que les deux autres as se trouvent ensemble dans la main d'un même joueur~?
\end{question}

\begin{question}\label{iq:iv1:timed-test-2:8}
Dix couples s'installent au hasard autour d'une table ronde. En moyenne, combien de couples sont-ils assis côte à côte~?
\end{question}

\begin{question}\label{iq:iv1:timed-test-2:9}
Un indice à pondérations égales contient 100 actions, chacune de volatilité 30~\%, et toute paire d'actions a la même corrélation $\rho$. L'indice a une volatilité de 15~\%. Quelle est la valeur de $\rho$~?
\end{question}

\begin{question}\label{iq:iv1:timed-test-2:10}
$X$ et $Y$ sont indépendantes et uniformes sur $[0, 1]$. Quelle est la probabilité que $X/Y$, arrondi à l'entier le plus proche, soit pair~?
\end{question}

\begin{question}\label{iq:iv1:timed-test-2:11}
Au tennis, un joueur remporte chaque point indépendamment avec la probabilité $0.6$. Un jeu est remporté en atteignant 4 points avec une avance d'au moins 2. Quelle est la probabilité que le joueur remporte le jeu~?
\end{question}

\begin{question}\label{iq:iv1:timed-test-2:12}
Calculez la somme $\sum_{n \ge 1}n^2/2^n$. Quel modèle probabiliste donne directement la réponse~?
\end{question}
